\begin{abstract}
Formula Student (FS) tracks are marked with hundreds of colour-coded cones that teams set out by hand with tapes and trigonometry, a slow process with no check of the result, although driverless cars depend on the layout matching the plan. This paper reviews the state of the art for TrackSeed, a system in which one person sets out a track with augmented reality (AR) guidance on a smartphone while a cloud-hosted digital twin (DT) follows the placement live. The review covers AR for on-site layout, mobile AR tracking and alignment to the ground, real-time synchronisation and on-device inference, using studies selected for quantitative evidence and official documentation for tool capabilities. The literature shows that AR placement reaches centimetre-to-decimetre accuracy over small areas, with phones the least accurate devices, and that tracking error grows with the distance from the reference; no study has evaluated phone-based setout over a full track. Based on these findings, TrackSeed adopts two-point marker alignment with re-alignment along the track, WebSocket synchronisation with a resend queue, dashboard edits sent back to the phone as guidance, and on-device cone verification and car identification. Placement accuracy, setup time, synchronisation latency and car identification accuracy are proposed to validate it.
\end{abstract}

\begin{IEEEkeywords}
augmented reality; digital twin; mobile AR tracking; spatial layout; Formula Student; real-time synchronisation; fine-grained vehicle classification
\end{IEEEkeywords}
